% Verdict table (code_council/build_verdicts.py); needs numbers_c_verdicts.tex.
% Like-for-like: code_council/peer_discrimination.py; strongest peer:
% code_council/strongest_peer.py (both Holm within setting).
\begin{table}[t]
\centering\small
\setlength\tabcolsep{4pt}
\begin{adjustbox}{max width=\linewidth}
\begin{tabular}{llllll}
\toprule
& & & \multicolumn{3}{c}{Self-specific advantage} \\
\cmidrule(lr){4-6}
Judges, texts, format & Self-naming test & Distinguishes own text & Pooled peers & Like-for-like & Strongest peer \\
\midrule
Frontier, chat-app essays, lineup & GPT, Claude & GPT, Claude & GPT, Claude & GPT, Claude & -- \\
Frontier, chat-app essays, one text & GPT, Claude & GPT, Claude & GPT & GPT, Claude & -- \\
Frontier, chat-app essays, yes/no & Claude & Claude & n/a & n/a & n/a \\
Frontier, API texts, lineup & Claude, Grok & Claude, Grok & -- & -- & -- \\
Frontier, API texts, one text & GPT, Claude & Claude & -- & Claude & -- \\
Open-weight, lineup & Qwen & -- & -- & -- & -- \\
Open-weight, one text & Qwen, Llama, Phi & -- & -- & -- & -- \\
Open-weight, yes/no & Llama, Mistral & -- & n/a & n/a & n/a \\
\midrule
Cases & \CVflagged{} of \CVcells{} & \CVdisc{} of \CVcells{} & \CVbothAttr{} of \CVattr{} & \CXLikeForLike{} of \CVattr{} & \CXStrongestPeer{} of \CVattr{} \\
\bottomrule
\end{tabular}
\end{adjustbox}
\caption{The audit applied to every case. It shows why the self-naming test
alone misleads: it flags \CVflagged{} cases, but in \CVfailFA{} of them the judge
names itself as often on other models' texts as on its own, and none of the
others is significantly better than its strongest peer. \emph{Self-naming test}: the judge
names itself on its own texts more often than chance. \emph{Distinguishes own
text}: it also names itself more often on its own texts than on other models'
texts (positive discrimination), so removing the author's name does not fully
hide authorship from it. \emph{Self-specific advantage}: it also identifies its
texts better than other judges do, measured three ways. \emph{Pooled peers}:
more often than the other judges name it on average (the peer baseline).
\emph{Like-for-like}: its discrimination exceeds the other judges' mean
discrimination for the same author. \emph{Strongest peer}: it names itself on
its texts more often than the single judge that names it most often there.
Every case in a later column also passes the self-naming test and
distinguishes its own text; the three peer columns are alternative measures,
not nested in one another. The strongest-peer test has little power with 38
essays per judge and is uninformative one text at a time, where peers name
Claude and GPT on nearly every one of their texts (Section~\ref{sec:setup}).
The yes/no format has no peer baseline (n/a), so the last three columns count
only the \CVattr{} lineup and one-text cases. Claude is the only yes/no case that distinguishes its own text, and it
does so only with correction within the panel ($p=$~\FBClaudeHolm{};
$p=\CXClaudeYesNoHolmAll$ across all \CVcells{} cases). All tests are
Holm-corrected across the judges of a panel.}
\label{tab:verdicts}
\end{table}
